\section{Introducción}
Abrir un recurso científico permite acceder a él; construir una comunidad exige organizar su uso, revisión y continuidad. Esta diferencia importa para instituciones venezolanas que buscan conectar formación, participación pública e investigación con recursos limitados. La Recomendación de la UNESCO sobre Ciencia Abierta sitúa la participación y la cooperación junto al acceso al conocimiento \cite{unesco}. El concepto de comunidad de práctica permite analizar el aprendizaje como participación sostenida en actividades compartidas \cite{wenger}.

Este ensayo sostiene que una actividad abierta se convierte en infraestructura comunitaria cuando deja tres elementos vinculados: una contribución identificable, un mecanismo de reconocimiento y una vía de continuación. La convocatoria y el número de asistentes describen el alcance de un encuentro, pero su continuidad se juega en las responsabilidades y recursos que quedan después. La revisión de Schulten y Chounta muestra que los estudios sobre aprendizaje en hackatones recurren con frecuencia a percepciones y carecen de procedimientos suficientemente estandarizados de evaluación \cite{hacklearning}. Por ello proponemos examinar productos y mecanismos concretos antes de atribuir efectos educativos.

Comparamos tres formatos: CoAfina, la edición colaborativa de conocimiento sanitario y la formación con datos abiertos de física. La contribución del ensayo es una matriz de decisiones que distingue lo que cada formato conserva, cómo reconoce el trabajo y qué necesita para prolongarlo. CoAfina es uno de los casos de una trayectoria más amplia de ciencia abierta en América Latina.

El mapeo de Costa, García-Holgado y Alvarez organiza marcos metodológicos, estructuras de organización y métricas de hackatones estudiantiles, y señala límites en su seguimiento longitudinal \cite{hackmapping}. Estas dimensiones fundamentan la separación entre organización, productos y continuidad en nuestra matriz.
